\paragraph{Notations.}
For nonnegative quantities \(a,b\), we write \(a\lesssim b\) if
\(a\le Cb\) for a constant \(C>0\) independent of the main asymptotic
parameters, such as the particle number \(N\), horizon \(T\), and
failure probability \(p\).  Equivalently, \(a=O(b)\) means \(a\lesssim b\).
We write \(a=o(b)\) if \(a/b\to0\) in the prescribed asymptotic regime.
Similarly, \(a=\Omega(b)\) means \(b=O(a)\), and \(a=\omega(b)\) means
\(b=o(a)\).  We write \(a=\Theta(b)\) if both \(a=O(b)\) and
\(a=\Omega(b)\) hold.  Finally, \(a=\widetilde O(b)\) hides
polylogarithmic factors in the main asymptotic parameters.  The hidden
constants may depend on fixed problem parameters, the dimension, and structural smoothness constants, but not on \(N,T\), or \(p\).

For normed spaces \(E,F\), \(\mathcal L(E;F)\) denotes the space of bounded
linear maps from \(E\) to \(F\), and \(\mathcal L(E_1,\ldots,E_k;F)\) denotes
the space of bounded \(k\)-linear maps from \(E_1\times\cdots\times E_k\) to
\(F\).  When \(E=F\), we write \(\mathcal L(E)\).  For a linear map
\(A:\mathbb R^d\to\mathbb R^d\), 
$
    \norm{A}_{\Loperator{d}}
    :=
    \sup_{\norm{v}_2=1}\norm{Av}_2
$
denotes its Euclidean operator norm.  We use \(\norm{\cdot}_{\HS}\) for the
Hilbert--Schmidt norm and \(\norm{\cdot}_{\F}\) for the Frobenius norm; in
finite-dimensional matrix spaces these coincide, but we reserve
\(\norm{\cdot}_{\HS}\) for Hilbert-space operators or tensor-valued feature
maps, and \(\norm{\cdot}_{\F}\) for ordinary finite matrices.

\paragraph{Vector fields, configurations, and empirical measures.}
For a measure \(\mu\in\Ptwo(\mathbb R^d)\) and vector fields
\(v,w:\mathbb R^d\to\mathbb R^d\), define
\[
 \|v\|_{L^2(\mu)}^2:=\int\|v(x)\|^2\,\mu(dx),
 \qquad
 \langle v,w\rangle_{L^2(\mu)}:=\int\langle v(x),w(x)\rangle\,\mu(dx).
\]
For a configuration \(\mathbf x=(x^1,\ldots,x^N)\in(\mathbb R^d)^N\), set
\[
 \|\mathbf x\|_N^2:=\frac1N\sum_{i=1}^N\|x^i\|^2,
 \qquad
 \rho_{\mathbf x}:=\frac1N\sum_{i=1}^N\delta_{x^i},
 \qquad
 \xi^N(\mathbf x):=(\xi(x^1),\ldots,\xi(x^N)).
\]
We abbreviate \(\rho_{X_t}:=\rho_{\mathbf X_t}\) and
\(\rho_{Y_t}:=\rho_{\mathbf Y_t}\). The coupled discrepancy is
\[
 \delta_{t,i}:=X_t^i-Y_t^i,
 \qquad
 r_t^2:=\frac1N\sum_{i=1}^N\|\delta_{t,i}\|^2
       =\|\mathbf X_t-\mathbf Y_t\|_N^2.
\]
The normalized configuration inner product is
\[
 \langle\mathbf u,\mathbf v\rangle_N
 :=\frac1N\sum_{i=1}^N\langle u^i,v^i\rangle.
\]

\paragraph{Gradient, Hessian, and jump covariance.}
The Wasserstein gradient is \(\dW F(\mu,x)\in\mathbb R^d\).
We distinguish its measure derivative from its spatial Jacobian:
\begin{align*}
 \HessW F(\mu;x,y)&:=\nabla_W^2F(\mu;x,y)\in\mathbb R^{d\times d},\\
 \xHessW F(\mu,x)&:=\nabla_x\nabla_WF(\mu,x)\in\mathbb R^{d\times d}.
\end{align*}
Use the following operators on their natural domains in
\(L^2(\mu;\mathbb R^d)\):
\begin{align*}
 (H_\mu v)(x)&:=\int\HessW F(\mu;x,y)v(y)\,\mu(dy),\\
 (H_\mu'v)(x)&:=\xHessW F(\mu,x)v(x),
 &\mathcal H_\mu&:=H_\mu+H_\mu'.
\end{align*}
In particular, the spatial multiplication operator has domain
\[
 \mathcal D(H_\mu')
 :=\{v\in L^2(\mu;\mathbb R^d):
       \xHessW F(\mu,\cdot)v(\cdot)\in L^2(\mu;\mathbb R^d)\},
 \qquad
 \mathcal D(\mathcal H_\mu)=\mathcal D(H_\mu)\cap\mathcal D(H_\mu').
\]
Here \(\mathcal D(H_\mu)\) consists of the fields for which the displayed
kernel integral exists and defines an \(L^2(\mu)\) field.
For $v\in\mathcal D(\mathcal H_\mu)$ for which the pushforward perturbation and
the indicated derivatives are defined, the full second-variation identity is
\[
 \frac{d^2}{dh^2}\Big|_{h=0}F((\Id+hv)_\#\mu)
 =\langle v,\mathcal H_\mu v\rangle_{L^2(\mu)}.
\]
The tracking curvature retains both terms of \(\mathcal H_\mu\).
The additional spatial-Hessian bound in
\cref{prop:finite-feature-negative-curvature} makes \(H_\mu'\) bounded;
no such global bound is imposed here.
In contrast, the Hessian-guided jump covariance uses the integral term:
\[
 K_\mu(x,y):=\int\HessW F(\mu;x,z)\HessW F(\mu;z,y)\,\mu(dz),
 \qquad \xi_\mu\sim\GP(0,K_\mu).
\]
Here \(\GP(0,K_\mu)\) denotes a centered Gaussian process with covariance
kernel \(K_\mu\). The finite Gaussian coordinates used to compare two
jumps are constructed from the population field in
\cref{lem:population-first-finite-coupling}.
For a matrix or trace-class operator \(A\), we use the nuclear norm
\[
 \|A\|_*:=\operatorname{tr}((A^*A)^{1/2}),
\]
where \(A^*\) denotes the Hilbert adjoint and equals \(A^\top\) for real matrices.

\paragraph{Conventions for the fixed-schedule comparison.}
The jump count is \(\NofJ\). We take
\(0<t_1<\cdots<t_{\NofJ}\le T\), put \(t_0=0\) and
\(t_{\NofJ+1}=T\), and allow \(\NofJ=0\).
Paths are right-continuous: \(Z_{t_\ell}\) is the post-jump value and
\(Z_{t_\ell-}=\lim_{s\uparrow t_\ell}Z_s\) is the pre-jump value.
Jump covariances and jump residuals are evaluated at \(t_\ell-\).
Empty sums and products are zero and one, respectively.

\paragraph{Population conditioning and comparison history.}
The population initial law and schedule are deterministic. Population
Gaussian innovations are independent of the iid initial sample.
Let \(\mathcal G_\ell\) be the sigma-field generated by the population
fields through jump \(\ell\), and let
\(\mathcal G_{\rm pop}=\mathcal G_{\NofJ}\).
The population path satisfies
\(\mu_t=\Law(Y_t^1\mid\mathcal G_{\rm pop})\);
the conditional product law is proved in
\cref{lemma-gp-jump-implies-conditional-iid}.
The notation \(\mathcal H_{\ell-}\) denotes the full pre-jump history,
including \(X_{t_\ell-},Y_{t_\ell-},\mu_{t_\ell-}\), all previous
population innovations and auxiliary normals, and the initial sample.
It does not include the current or future innovations. Finite Gaussian comparison
coordinates are standard Gaussian conditionally on this history;
they are not the conditioning object in the product law for \(Y\).
